% Introduction for Malleable World Models (MWM)
% % Introduction for Malleable World Models (MWM)
% % Introduction for Malleable World Models (MWM)
% % Introduction for Malleable World Models (MWM)
% \input{introduction}

\section{Introduction}
\label{sec:introduction}

World models enable model-based agents to plan and act by predicting latent dynamics from experience~\cite{hafner2023mastering,hafner2020dream}.
Once trained, however, they are typically deployed with \emph{frozen} parameters.
In open-ended settings---wear on actuators, payload changes, or simulator-to-real gaps---the transition kernel internalized during pretraining ceases to match the environment the agent inhabits.
Prediction error rises, imagined rollouts become misleading, and control performance degrades.
\emph{Test-time adaptation} (TTA) seeks to recover competence after such shifts without returning to expensive offline training.

The dominant TTA recipe is gradient-based learning at deployment: collect post-shift trajectories, backpropagate a surrogate loss, and update network weights~\cite{nagabandi2018learning,zintgraf2021fast}.
This \emph{slow} pathway is expressive but incurs costs that matter in interactive control---batched updates, autograd through long horizons, and tuning to avoid catastrophic forgetting.
When the shift is sudden, the agent may act suboptimally for many steps before the first weight update.

Biological systems often separate a fast, local error-driven store from slower structural learning~\cite{mcclelland1995complementary}.
The fast store is not globally optimal; it provides coarse correction so behavior is not blind during transients.
We transpose this idea to world models: adaptation is not only a motor reflex, but a change in \emph{what the agent believes will happen next}.

We propose \textbf{Malleable World Models (MWM)}: frozen pretrained agents augmented with a \textbf{fast plastic prior residual} $P_t$, updated every step by a Hebbian rule without backprop through pretraining weights, plus optional \textbf{slow} actor and world-model updates from replay.
We evaluate on proprioceptive \texttt{walker-walk} under actuator gain shift ($1.0 \rightarrow 0.7 \rightarrow 1.0$).
Section~\ref{sec:method} details the method; Section~\ref{sec:experiments} reports detection, adaptation latency, and consolidation behavior.

\paragraph{Contributions.}
\begin{enumerate}[leftmargin=*, itemsep=2pt]
  \item A backpropagation-free plastic prior on RSSM logits with norm-capped state $P_t$.
  \item A tri-timescale test-time stack: fast $P_t$, slow actor fine-tuning, and optional WM consolidation with coupled flush.
  \item Empirical characterization on gain shift: rapid trace growth and prior-error changes; control gains when fast plastic scaffolds slow actor learning; pilot evidence for $P_t$ drainage under consolidation.
\end{enumerate}


\section{Introduction}
\label{sec:introduction}

World models enable model-based agents to plan and act by predicting latent dynamics from experience~\cite{hafner2023mastering,hafner2020dream}.
Once trained, however, they are typically deployed with \emph{frozen} parameters.
In open-ended settings---wear on actuators, payload changes, or simulator-to-real gaps---the transition kernel internalized during pretraining ceases to match the environment the agent inhabits.
Prediction error rises, imagined rollouts become misleading, and control performance degrades.
\emph{Test-time adaptation} (TTA) seeks to recover competence after such shifts without returning to expensive offline training.

The dominant TTA recipe is gradient-based learning at deployment: collect post-shift trajectories, backpropagate a surrogate loss, and update network weights~\cite{nagabandi2018learning,zintgraf2021fast}.
This \emph{slow} pathway is expressive but incurs costs that matter in interactive control---batched updates, autograd through long horizons, and tuning to avoid catastrophic forgetting.
When the shift is sudden, the agent may act suboptimally for many steps before the first weight update.

Biological systems often separate a fast, local error-driven store from slower structural learning~\cite{mcclelland1995complementary}.
The fast store is not globally optimal; it provides coarse correction so behavior is not blind during transients.
We transpose this idea to world models: adaptation is not only a motor reflex, but a change in \emph{what the agent believes will happen next}.

We propose \textbf{Malleable World Models (MWM)}: frozen pretrained agents augmented with a \textbf{fast plastic prior residual} $P_t$, updated every step by a Hebbian rule without backprop through pretraining weights, plus optional \textbf{slow} actor and world-model updates from replay.
We evaluate on proprioceptive \texttt{walker-walk} under actuator gain shift ($1.0 \rightarrow 0.7 \rightarrow 1.0$).
Section~\ref{sec:method} details the method; Section~\ref{sec:experiments} reports detection, adaptation latency, and consolidation behavior.

\paragraph{Contributions.}
\begin{enumerate}[leftmargin=*, itemsep=2pt]
  \item A backpropagation-free plastic prior on RSSM logits with norm-capped state $P_t$.
  \item A tri-timescale test-time stack: fast $P_t$, slow actor fine-tuning, and optional WM consolidation with coupled flush.
  \item Empirical characterization on gain shift: rapid trace growth and prior-error changes; control gains when fast plastic scaffolds slow actor learning; pilot evidence for $P_t$ drainage under consolidation.
\end{enumerate}


\section{Introduction}
\label{sec:introduction}

World models enable model-based agents to plan and act by predicting latent dynamics from experience~\cite{hafner2023mastering,hafner2020dream}.
Once trained, however, they are typically deployed with \emph{frozen} parameters.
In open-ended settings---wear on actuators, payload changes, or simulator-to-real gaps---the transition kernel internalized during pretraining ceases to match the environment the agent inhabits.
Prediction error rises, imagined rollouts become misleading, and control performance degrades.
\emph{Test-time adaptation} (TTA) seeks to recover competence after such shifts without returning to expensive offline training.

The dominant TTA recipe is gradient-based learning at deployment: collect post-shift trajectories, backpropagate a surrogate loss, and update network weights~\cite{nagabandi2018learning,zintgraf2021fast}.
This \emph{slow} pathway is expressive but incurs costs that matter in interactive control---batched updates, autograd through long horizons, and tuning to avoid catastrophic forgetting.
When the shift is sudden, the agent may act suboptimally for many steps before the first weight update.

Biological systems often separate a fast, local error-driven store from slower structural learning~\cite{mcclelland1995complementary}.
The fast store is not globally optimal; it provides coarse correction so behavior is not blind during transients.
We transpose this idea to world models: adaptation is not only a motor reflex, but a change in \emph{what the agent believes will happen next}.

We propose \textbf{Malleable World Models (MWM)}: frozen pretrained agents augmented with a \textbf{fast plastic prior residual} $P_t$, updated every step by a Hebbian rule without backprop through pretraining weights, plus optional \textbf{slow} actor and world-model updates from replay.
We evaluate on proprioceptive \texttt{walker-walk} under actuator gain shift ($1.0 \rightarrow 0.7 \rightarrow 1.0$).
Section~\ref{sec:method} details the method; Section~\ref{sec:experiments} reports detection, adaptation latency, and consolidation behavior.

\paragraph{Contributions.}
\begin{enumerate}[leftmargin=*, itemsep=2pt]
  \item A backpropagation-free plastic prior on RSSM logits with norm-capped state $P_t$.
  \item A tri-timescale test-time stack: fast $P_t$, slow actor fine-tuning, and optional WM consolidation with coupled flush.
  \item Empirical characterization on gain shift: rapid trace growth and prior-error changes; control gains when fast plastic scaffolds slow actor learning; pilot evidence for $P_t$ drainage under consolidation.
\end{enumerate}


\section{Introduction}
\label{sec:introduction}

World models enable model-based agents to plan and act by predicting latent dynamics from experience~\cite{hafner2023mastering,hafner2020dream}.
Once trained, however, they are typically deployed with \emph{frozen} parameters.
In open-ended settings---wear on actuators, payload changes, or simulator-to-real gaps---the transition kernel internalized during pretraining ceases to match the environment the agent inhabits.
Prediction error rises, imagined rollouts become misleading, and control performance degrades.
\emph{Test-time adaptation} (TTA) seeks to recover competence after such shifts without returning to expensive offline training.

The dominant TTA recipe is gradient-based learning at deployment: collect post-shift trajectories, backpropagate a surrogate loss, and update network weights~\cite{nagabandi2018learning,zintgraf2021fast}.
This \emph{slow} pathway is expressive but incurs costs that matter in interactive control---batched updates, autograd through long horizons, and tuning to avoid catastrophic forgetting.
When the shift is sudden, the agent may act suboptimally for many steps before the first weight update.

Biological systems often separate a fast, local error-driven store from slower structural learning~\cite{mcclelland1995complementary}.
The fast store is not globally optimal; it provides coarse correction so behavior is not blind during transients.
We transpose this idea to world models: adaptation is not only a motor reflex, but a change in \emph{what the agent believes will happen next}.

We propose \textbf{Malleable World Models (MWM)}: frozen pretrained agents augmented with a \textbf{fast plastic prior residual} $P_t$, updated every step by a Hebbian rule without backprop through pretraining weights, plus optional \textbf{slow} actor and world-model updates from replay.
We evaluate on proprioceptive \texttt{walker-walk} under actuator gain shift ($1.0 \rightarrow 0.7 \rightarrow 1.0$).
Section~\ref{sec:method} details the method; Section~\ref{sec:experiments} reports detection, adaptation latency, and consolidation behavior.

\paragraph{Contributions.}
\begin{enumerate}[leftmargin=*, itemsep=2pt]
  \item A backpropagation-free plastic prior on RSSM logits with norm-capped state $P_t$.
  \item A tri-timescale test-time stack: fast $P_t$, slow actor fine-tuning, and optional WM consolidation with coupled flush.
  \item Empirical characterization on gain shift: rapid trace growth and prior-error changes; control gains when fast plastic scaffolds slow actor learning; pilot evidence for $P_t$ drainage under consolidation.
\end{enumerate}
